\documentclass[11pt]{article}
\usepackage[margin=1in]{geometry}
\usepackage{amsmath,amssymb,amsthm}
\usepackage[T1]{fontenc}
\usepackage{lmodern}
\usepackage{graphicx}
\usepackage[hidelinks]{hyperref}
\usepackage{xcolor}

\newtheorem{theorem}{Theorem}
\newtheorem{lemma}[theorem]{Lemma}
\newcommand{\SG}{\widetilde{SG}}
\newcommand{\J}{\mathcal J}
\newcommand{\Rmap}{R}

\title{An Extra Residual Spectral Point on the\\
Barlow--Perkins Sierpinski Lattice}
\author{Candidate counterexample for arXiv:1910.01771}
\date{11 August 2026}

\begin{document}
\maketitle

\begin{center}
\fcolorbox{black}{gray!8}{\parbox{0.88\linewidth}{
\textbf{Status: candidate counterexample, likely valid.}
The Barlow--Perkins question is answered negatively, subject to human review.}}
\end{center}

\section{Source question}

Shiping Cao, Yiqi Huang, Hua Qiu, Robert S. Strichartz, and Xiaohan Zhu
study the graph Laplacian on infinite Sierpinski lattices in
\cite{CHQSZ}. In Section 3.3 (arXiv:1910.01771, PDF page 14), they write
the Barlow--Perkins lattice $\SG_{BP}$ as the union of two boundary-point
Sierpinski half-lattices and observe that reflection splits its Laplacian into
a Neumann and a Dirichlet half. After proving the $1<p<\infty$ spectral
classification, they ask:

\begin{quote}
On the other hand, the argument in Proposition 3.4, showing that the
eigenfunctions are not bounded, relies heavily on (3.2). It is not clear how
to overcome this, and we do not know whether the same result of Theorem 3.3
(b) holds on $\SG_{BP}$.
\end{quote}

\clearpage
\begin{figure}[p]
  \centering
  \includegraphics[width=0.9\textwidth]{figures/open_problem_crop.png}
  \caption{The Barlow--Perkins discussion and question in
  arXiv:1910.01771, PDF page 14.}
\end{figure}
\clearpage

Theorem 3.3(b) says, for a one-sided lattice with a boundary point, that the
$\ell^1$ residual spectrum is $\{0\}\cup\Sigma_4$ and that
\[
 \J\setminus(\Sigma_4\cup\Sigma_5\cup\{0\})
\]
is continuous spectrum. Here
\[
 \Rmap(t)=t(5-t),\qquad
 \Sigma_4=\bigcup_{m\ge0}\Rmap^{-m}\{4\},\qquad
 \Sigma_5=\bigcup_{m\ge0}\Rmap^{-m}\{5\},
\]
and
\[
 \Sigma_6=\{6\}\cup\bigcup_{m\ge0}\Rmap^{-m}\{3\}.
\]
As in the source, a ``$\lambda$-eigenfunction'' satisfies
$-\Delta f=\lambda f$.

\section{Proof intuition}

The missing odd-reflection half is a Dirichlet problem. We choose the
period-two point
\[
 a=3-\sqrt3,\qquad \Rmap(a)=3+\sqrt3=:b,\qquad \Rmap(b)=a.
\]
Two steps of spectral decimation therefore repeat exactly. Symmetric boundary
data on the Dirichlet half are multiplied by $-3$ when transported two levels
inward, so prescribing outer data $(-1/3)^n(0,1,1)$ on a level-$2n$ gasket
makes all finite approximants consistent.

Consistency alone does not prove boundedness: individual extension
coefficients are larger than one. The key certificate is a positive quadratic
form $P$ for which every one of the nine possible two-level child-extension
matrices has norm at most $3$. The prefactor $3^{-n}$ cancels this worst-case
growth exactly. The consistent finite solutions therefore converge to a
nonzero bounded Dirichlet eigenfunction. Odd reflection across the gluing
vertex gives a bounded $a$-eigenfunction on $\SG_{BP}$. Duality then puts $a$
in the $\ell^1$ residual spectrum, contrary to the proposed extension of
Theorem 3.3(b).

\section{The two-level extension certificate}

For a finite gasket cell, order the three boundary values as
$v=(v_0,v_1,v_2)^T$. Proposition 2.3 of \cite{CHQSZ} gives the usual spectral
decimation extension rule. If the fine eigenvalue is $t$, the value at the
midpoint of the edge joining vertices $i,j$ is
\[
 \frac{(4-t)(v_i+v_j)+2v_k}{(2-t)(5-t)},                     \tag{1}
\]
where $\{i,j,k\}=\{0,1,2\}$.

For $t=a$ and $t=b$, the coefficient of $v_i+v_j$ in (1) is one, while the
coefficient of $v_k$ is respectively
\[
 c_a=\sqrt3-1,\qquad c_b=-(1+\sqrt3).                        \tag{2}
\]
For a scalar $c$, define the three child-boundary matrices
\[
 L_0(c)=\begin{pmatrix}1&0&0\\1&1&c\\1&c&1\end{pmatrix},\quad
 L_1(c)=\begin{pmatrix}1&1&c\\0&1&0\\c&1&1\end{pmatrix},\quad
 L_2(c)=\begin{pmatrix}1&c&1\\c&1&1\\0&0&1\end{pmatrix}. \tag{3}
\]
Thus a two-level child has boundary vector
\[
 M_{ij}v,\qquad M_{ij}=L_j(c_a)L_i(c_b),\qquad 0\le i,j\le2. \tag{4}
\]

\begin{lemma}[Common quadratic bound]\label{lem:matrix}
Let $J$ denote the $3\times3$ all-ones matrix and put
\[
 P=I-\frac{\sqrt3}{6}J.
\]
Then $P$ is positive definite and, for all $0\le i,j\le2$,
\[
 M_{ij}^{T}PM_{ij}\le9P.                                    \tag{5}
\]
Moreover, for $v_*=(0,1,1)^T$,
\[
 M_{00}v_*=-3v_*.                                            \tag{6}
\]
\end{lemma}

\begin{proof}
The eigenvalues of $P$ are $1,1,1-\sqrt3/2$, so $P>0$. The matrices in
(4) are equivariant under simultaneous permutations of the three
coordinates, and $P$ is permutation invariant. Hence there are only two
cases: $i=j$ and $i\ne j$.

For representatives $M_{00}$ and $M_{01}$, direct multiplication gives
\[
 9P-M_{00}^{T}PM_{00}=
 \begin{pmatrix}
 6-2\sqrt3&0&0\\0&4&-4\\0&-4&4
 \end{pmatrix}\ge0,                                        \tag{7}
\]
and
\[
 9P-M_{01}^{T}PM_{01}=
 \begin{pmatrix}
 6-2\sqrt3&(3-\sqrt3)/2&-\sqrt3\\
 (3-\sqrt3)/2&4-\sqrt3&(-5+\sqrt3)/2\\
 -\sqrt3&(-5+\sqrt3)/2&4
 \end{pmatrix}.                                             \tag{8}
\]
The leading principal minors of (8) are
\[
 6-2\sqrt3,\qquad \frac{54-25\sqrt3}{2},\qquad
 9(3-\sqrt3),
\]
all strictly positive. Thus (8) is positive definite and (5) follows.
Finally, (3)--(4) give (6) directly.
\end{proof}

Write $\|v\|_P=(v^TPv)^{1/2}$. Lemma~\ref{lem:matrix} says
\[
 \|M_{ij}v\|_P\le3\|v\|_P.                                 \tag{9}
\]

\section{Construction of the bounded eigenfunction}

\begin{theorem}\label{thm:counterexample}
Let $a=3-\sqrt3$. Then $a$ is an $\ell^\infty$ eigenvalue of the
Barlow--Perkins Sierpinski lattice Laplacian. Consequently, $a$ belongs to
the residual spectrum of the Laplacian on $\ell^1(\SG_{BP})$.
In particular, the analogue of Theorem 3.3(b) in \cite{CHQSZ} is false on
$\SG_{BP}$.
\end{theorem}

\begin{proof}
We first work on one boundary-point half with Dirichlet value zero at its
boundary vertex. Let $G_{2n}$ be the nested level-$2n$ finite gasket in that
half. At its three outer corners prescribe
\[
 d_n v_*,\qquad d_n=\left(-\frac13\right)^n,\qquad
 v_*=(0,1,1)^T.                                              \tag{10}
\]
Starting with this outer boundary vector, extend through $n$ pairs of levels,
first with the $b$ rule and then with the $a$ rule in each pair. Formula (1)
shows that this produces a function $f_n$ satisfying
$-\Delta f_n=af_n$ at every vertex of $G_{2n}$ except its three outer
corners.

The distinguished child obtained by choosing the $0$-cell twice has boundary
vector
\[
 d_nM_{00}v_*=-3d_nv_*=d_{n-1}v_*
\]
by (6). Repeating inside that child shows that $f_n$ restricts to $f_{n-1}$
on the naturally embedded $G_{2n-2}$. The finite functions are therefore
consistent.

Every vertex value of $f_n$ is a coordinate of
\[
 d_nM_{i_nj_n}\cdots M_{i_1j_1}v_*                         \tag{11}
\]
for some sequence of $n$ two-level children. By (9), the $P$-norm of (11)
is at most
\[
 |d_n|3^n\|v_*\|_P=\|v_*\|_P.                              \tag{12}
\]
Since $P>0$, Euclidean coordinate norm is bounded by a fixed multiple of
$\|\cdot\|_P$. Thus all $f_n$ are bounded by one constant independent of
$n$.

The consistent union $f_D=\bigcup_nf_n$ is consequently a bounded function
on the infinite Dirichlet half. Every nonboundary vertex is an interior
vertex of some $G_{2n}$, so $-\Delta_Df_D=af_D$ everywhere. It is nonzero:
on the two non-Dirichlet vertices of the initial cell its values are both one.

Now use the reflection of $\SG_{BP}$ exchanging its two halves. Put $f_D$ on
one half, $-f_D$ on the other, and zero at the common vertex. The eigenvalue
equation holds away from the common vertex by construction. At the common
vertex, the four neighboring values cancel in reflected pairs, so the
equation holds there as well. This gives a nonzero bounded $a$-eigenfunction
on $\SG_{BP}$.

The orbit $a\mapsto b\mapsto a$ is bounded, hence $a\in\J$. It never meets
$0,3,4,5$, or $6$, so
\[
 a\notin\{0\}\cup\Sigma_4\cup\Sigma_5\cup\Sigma_6.       \tag{13}
\]
In particular, Theorem 3.10 of \cite{CHQSZ} at $p=2$ shows that $a$ is not
an $\ell^2$ eigenvalue; since $\ell^1\subset\ell^2$, it is not an
$\ell^1$ eigenvalue either.

On the other hand, the bounded eigenfunction just constructed belongs to the
kernel of the adjoint of $a+\Delta$ acting on $\ell^1$. Hence the range of
$a+\Delta$ on $\ell^1$ is not dense. Since $a+\Delta$ is injective there,
$a$ is a residual spectral point. Theorem 3.3(b) would instead place (13) in
the continuous spectrum, proving the claimed negative answer.
\end{proof}

\section{Verification, limitations, and novelty check}

The exact matrix certificate is independently reproduced by
\path{code/verify_matrix_certificate.py}. It checks all nine inequalities
in (5), positivity of $P$, and (6) symbolically over
$\mathbb Q(\sqrt3)$. The auxiliary
\path{code/finite_level_probe.py} builds finite gasket graphs and solves
the boundary problems numerically. Levels $1$ through $8$ (up to $9843$
vertices) retain the inner values $(1,1)$ and have observed supremum
$2/\sqrt3$. These computations are checks, not ingredients of the proof.

This counterexample fully answers the specific yes/no Barlow--Perkins
question negatively. It does not classify the entire $\ell^1$ spectrum of
$\SG_{BP}$, nor does it answer the broader preceding problem about all
$\ell^\infty$ eigenvalues of arbitrary boundaryless Sierpinski lattices or
their dependence on the generating word.

A bounded literature search checked arXiv:1910.01771, the exact sentence
``we do not know whether the same result'', and combinations of
``Barlow--Perkins'', ``Sierpinski lattice'', ``$\ell^1$ spectrum'',
``$\ell^\infty$ eigenvalue'', ``generating sequence'', and
``$3-\sqrt3$''. Searches also checked later arXiv results on Sierpinski
lattice spectra and spectral dimension. No later source answering this
question or giving the eigenvalue above was located. This supports, but does
not certify, novelty.

\paragraph{Human-review recommendation.}
Review as a likely valid full negative answer to the Barlow--Perkins question.
The highest-value checks are the orientation/order of the two spectral
decimation extensions in (4), consistency of the nested finite gaskets, and
the final $\ell^\infty$--$\ell^1$ adjoint argument.

\begin{thebibliography}{9}

\bibitem{CHQSZ}
S. Cao, Y. Huang, H. Qiu, R. S. Strichartz, and X. Zhu,
\emph{Spectral analysis beyond $\ell^2$ on Sierpinski lattices},
arXiv:1910.01771, version 2 (2021),
\url{https://arxiv.org/abs/1910.01771}.

\end{thebibliography}

\end{document}
